\begin{lemma}[Exponentially accurate aggregate geometry above a logarithmic scale]
\label{restore:high-aggregate}
Fix any $A_0>0$ and $H>0$. Under the fixed positive compact structural
conditions, if $d\ge A_0\log n$ and $np^2\le C_1$, there is an aggregate
event $\mathcal E_H$ with
\[
 \PP(\mathcal E_H^c)\le C_H e^{-Hd}
\]
on which the empirical-floor geometry, weighted degree-energy bounds,
two-sided total-degree bounds, and true block-total bounds used in
\cref{restore:fit-geometry-one,restore:fit-geometry-two,restore:fit-block-event}
hold for the full graph and every graph obtained by deleting at most three
vertices. The constants may depend on $A_0,H,C_1$ and the structural
bounds. No condition is imposed on the smallest raw degree or the
largest spectral row.
\end{lemma}

\begin{proof}
The centered bounded-entry norm inequality of Bandeira--van Handel
\cite[Corollary 3.12 and Remark 3.13]{bvh}, as used in
\eqref{eq:raw-norm-tail}, gives
\[
 \PP\{\norm{A-\Omega}_{\op}>L_H\sqrt d\}
 \le e^{-(H+2)d}
\]
after choosing a sufficiently large fixed $L_H$. Indeed, the row
variance sum is $O(d)$, the entry bound is one, and the prefactor $n$
in the norm-tail inequality is absorbed using $\log n\le d/A_0$.
The completed-mean diagonal has norm $O(p)$ and is absorbed in the
same threshold. This step is solely a norm event; it is not intersected
with the raw minimum-degree event in \eqref{eq:base-event}.

The Bernoulli transform of the total number of edges gives, with
failure probability at most $Ce^{-cnd}$,
\[
 cnd\le\sum_iD_i\le Cnd.
\]
For each of the fixed $K^2$ true block totals, Bernstein's inequality
for the independent upper-triangular edges, whose coefficients have
absolute value at most two, gives
\[
 \PP\left\{\left|[Z^\top(A-\Omega)Z]_{ab}\right|>Ln\right\}
 \le 2\exp\left\{-\frac{cL^2n}{d+L}\right\}.
\]
The deterministic completed diagonal costs $O(d)$ and is included
by increasing $L$. Since $d\to\infty$ and $d^2/n\le C_1$,
the exponent is at least $c'L^2d/C_1$ for all sufficiently large $n$.
Increasing the fixed $L=L_H$ therefore makes the failure of all block
totals at most $C_He^{-(H+2)d}$. The total-degree failure is also
bounded by $Ce^{-(H+2)d}$ for large $n$.

These three full-graph events imply the required events on every
fixed-depth minor without a degree-wise probability union. To see this,
let $J=[n]\setminus T$ with $|T|\le3$, set $A_J=A_{JJ}$ and
$\Omega_J=\Omega_{JJ}$, and recompute all degrees and the empirical
floor on $J$. The centered matrix $A_J-\Omega_J$ is a principal
submatrix of $A-\Omega$, so its norm is at most $L_H\sqrt d$.
Removing $T$ changes the total raw degree by at most $2|T|n$;
as $d\to\infty$, the two-sided bound by fixed multiples of $nd$
therefore persists uniformly over $T$. Removing $T$ changes a true
block total of $A-\Omega$ by $O(|T|n)$ deterministically, since
the entries have bounded absolute value. Its $O_H(n)$ bound persists
as well. The surviving group masses differ from their original
values by $O(1)$ and remain comparable to $n$.

The recomputed average degree is consequently comparable to $d$.
For the fixed sufficiently small floor multiplier $a_0$, its threshold
$\tau_J=a_0\overline D_J$ lies between positive fixed multiples of
$d$ and below every coordinate of $D_{0,J}=\Omega_J\one$.
Thus $G_{J,j}=\max(D_{J,j},\tau_J)\ge cd$ and
$\sum_jG_{J,j}\le Cnd$. The raw norm bound immediately gives
\begin{align*}
 \norm{G_J^{-1/2}(A_J-\Omega_J)G_J^{-1/2}}_{\op}
   &\le C_Hd^{-1/2},\\
 \sum_{j\in J}\frac{(D_{J,j}-D_{0,J,j})^2}{G_{J,j}}
   &=\norm{G_J^{-1/2}(A_J-\Omega_J)\one}_2^2
     \le C_Hn.
\end{align*}
Since the population floor is inactive,
$|G_{J,j}-D_{0,J,j}|\le|D_{J,j}-D_{0,J,j}|$; this proves the
other weighted energy bound. The fixed signal gap and the residual
$\norm{(I-U_JU_J^\top)G_J^{1/2}Z_J}_F\le C_H\sqrt n$
now follow by the deterministic Gram-matrix and projection arguments
in \cref{prop:sparse-floor}. Those arguments use only the floor,
total degree, normalized noise, and degree energy, all just proved.
Intersecting the stated aggregate consequences defines
$\mathcal E_H$. The three full-graph events imply it simultaneously
for every indicated minor, proving its claimed failure bound.

For a deleted graph used as independent training data, form the
analogous event using that graph and its further minors. The same
argument gives failure probability $C_He^{-Hd}$, and this event is
measurable from the training graph alone. The full-graph common event
is used only for deterministic implications, never as a condition
that would preserve independence of an incident row.
\end{proof}

\begin{proof}[Completion over the remaining degree scales]
Fix $A_0>0$ and use the preceding sparse-range argument when
$d\le A_0\log n$. For $d>A_0\log n$, replace its polynomially
accurate aggregate event by $\mathcal E_H$ from
\cref{restore:high-aggregate}. The proof of
\cref{restore:tail-energy-lemma} uses only the deterministic aggregate
geometry on the event; the innovation concentration is applied for
every training graph before restricting to the event. It therefore
holds unchanged with $\mathcal E_H$ and the corresponding
training-measurable events. Its coefficient-hit and target-row
corollaries, the scalar LOO comparison, all fitted-score bounds,
and the invariant basin likewise apply.

The resulting risk upper bound is the appropriate oracle term plus
$C_{H'}e^{-H'd}+C_He^{-Hd}$, where $H'$ can be prescribed by choosing
the fixed local thresholds. On either fixed positive compact model
class the oracle risk is at least $cd^{-1/2}e^{-C_Id}$.
Choosing $H,H'>C_I+1$ pays both remainders. This proves the stated
all-scale bound without a raw information--degree margin or an
entry-ratio-two restriction for the regularized ordinary block-LRT
decoder. The ordinary proof continues to retain its full Bernoulli
nonedge intercept, and the DC proof continues to compare first to
the profile oracle. The known-target-degree DC comparison retains
its separate equal-exposure assumption.
\end{proof}
